\section{Puncturing, shortening, and relative recovery costs}
\label{sec:relative-weights}

The local problem forgets the outer code and asks how many helpers realize
a prescribed target combination.  In the running example, the helper
coefficient vectors $(1,0,0,0)$, $(0,1,0,0)$, and $(0,0,1,1)$ all realize
$u$.  Their differences lie in the kernel of the helper generator map.
Passing to its quotient identifies the requested combination while leaving
the cost of its different physical realizations to be minimized.

Continue with $G=(G_P\mid G_J)$, viewed as
$G_P:\F_q^P\to\F_q^k$ and $G_J:\F_q^J\to\F_q^k$, and set
\[
 U_P=\operatorname{im}G_P,
 \qquad
 W_P=U_P\cap\operatorname{im}G_J,
\]
\[
 K_P=\ker G_J,
 \qquad
 D_P=G_J^{-1}(U_P).
\]
\begin{proposition}[Associated nested code pair]
\label{prop:associated-pair}
\coverage{complete}%
\lean{helper_ker_le_helperCodeForTargetSpace,%
map_helperCodeForTargetSpace_eq_recoverableTargetMessageSpace,%
recoverableTargetMap_surjective,%
mem_ker_recoverableTargetMap_iff}%
Restriction of $G_J$ gives an exact sequence
\begin{equation}
 0\longrightarrow K_P\longrightarrow D_P
 \xrightarrow{G_J}W_P\longrightarrow0.
 \label{eq:associated-exact-sequence}
\end{equation}
\end{proposition}

\begin{proof}
By definition, $K_P=\ker G_J$ is contained in
$D_P=G_J^{-1}(U_P)$.  The kernel of the restricted map is therefore $K_P$.
Its image is
\[
 G_J(D_P)=\operatorname{im}G_J\cap U_P=W_P,
\]
so the restricted map is onto $W_P$.
\proves{prop:associated-pair}%
\end{proof}

The first isomorphism theorem applied to the restricted map gives
$D_P/K_P\cong W_P$.

\begin{proposition}[Canonical puncturing--shortening description]
\label{prop:puncture-shorten-pair}
\coverage{absent}%
With puncturing and shortening taken onto the helper set $J$,
\begin{equation}
 D_P=\operatorname{punct}_J(I^\perp),
 \qquad
 K_P=\operatorname{short}_J(I^\perp).
 \label{eq:puncture-shorten-pair}
\end{equation}
Thus $\operatorname{punct}_J$ restricts to $J$, while
$\operatorname{short}_J$ first requires the coordinates on $P$ to vanish and
then restricts to $J$.  These operations are dual to one another.  Taking
duals inside $\F_q^J$ gives
\begin{equation}
 D_P^\perp=\operatorname{short}_J(I),
 \qquad
 K_P^\perp=\operatorname{punct}_J(I).
 \label{eq:dual-puncture-shorten-pair}
\end{equation}
Consequently the nested pair depends only on the code $I$ and the split
$E=P\sqcup J$; changing the row basis of a generator matrix leaves it
unchanged.
\end{proposition}

\begin{proof}
A vector $a\in\F_q^J$ lies in $\operatorname{punct}_J(I^\perp)$ exactly when
some $p\in\F_q^P$ satisfies $G_Pp+G_Ja=0$, equivalently
$G_Ja\in\operatorname{im}G_P$, which is $a\in D_P$.  It lies in
$\operatorname{short}_J(I^\perp)$ exactly when $(0,a)\in I^\perp$,
equivalently $G_Ja=0$, which is $a\in K_P$.
The dual identities follow from the standard puncturing--shortening duality.
\proves{prop:puncture-shorten-pair}%
\end{proof}

Thus a monomial relabeling of helpers gives the corresponding monomially
equivalent pair, while a row-basis change does nothing to the pair itself.

For a subspace $A\leq\F_q^J$, write $\supp A$ for the union of the supports
of its vectors.  For nested linear codes $K\subseteq D\leq\F_q^J$, the $t$th relative
generalized Hamming weight is
\[
 M_t(D,K)=
 \min\{ |\supp L|:L\leq D,
     \dim L-\dim(L\cap K)=t\}.
\]
Equivalently, one may minimize over subspaces disjoint from $K$ and of
dimension $t$.  Strict growth and the relative Singleton bound are
Proposition~2 and Theorem~4 of Luo et al.~\cite{LuoEtAl2005}; standard
treatments also cover ordinary Wei duality and computational forms
\cite{GeilEtAl2014,SanJose2025}.
For nested code pairs in linear secret sharing, RDLPs and RGHWs already give
exact information and reconstruction thresholds
\cite{KuriharaUyematsuMatsumoto2012}.

The following result is the standard nested-code support invariant applied to
the canonical puncturing--shortening pair, with the normalization and helper
support correspondence stated in recovery language.

\begin{theorem}[Relative weights are exact recovery costs]
\label{thm:relative-weight-recovery}
\coverage{absent}%
Let $\ell=\dim W_P$ and $1\leq t\leq\ell$.  The minimum helper-union size
among recovery systems for all helper-recoverable $t$-dimensional
subspaces of target linear combinations, denoted $\mu_t(P)$, is
\[
 \mu_t(P)=M_t(D_P,K_P).
\]
More precisely, after forgetting the target normalization, normalized recovery
systems of recovered dimension $t$ are in support-preserving correspondence
with subspaces
\[
 L\leq D_P,\qquad \dim L=t,\qquad L\cap K_P=0,
\]
while normalization on the chosen recovered subspace $T\leq W_P$ is a linear
section of $G_P$ over $T$.
\end{theorem}

\begin{proof}
Let a normalized system recover a $t$-dimensional subspace $T\leq W_P$.
In the notation of Section~\ref{sec:recovery-equations}, put $L=\beta(T)$.
Since $G_J\beta=-\operatorname{id}_T$, the map $\beta$ is injective,
$L\leq D_P$, and $L\cap K_P=0$.  Moreover $G_JL=T$, and the union of
helpers in the system is exactly $\supp L$.

Conversely, let $L\leq D_P$ have dimension $t$ and meet $K_P$ trivially.
Then $G_J|_L$ is an isomorphism from $L$ onto the $t$-dimensional subspace
$T=G_JL\leq W_P$.  Choose a linear section
$\alpha:T\to\F_q^P$ with $G_P\alpha=\operatorname{id}_T$.  Put
$\beta=-(G_J|_L)^{-1}$.  Then
$G_P\alpha+G_J\beta=0$, so $(\alpha,\beta)$ is a normalized recovery system
with helper union $\supp L$.

It remains to compare this minimum with the standard relative-weight
definition.  If $L'\leq D_P$ satisfies
\[
 \dim L'-\dim(L'\cap K_P)=t,
\]
choose a vector-space complement $L$ to $L'\cap K_P$ inside $L'$.  Then
$\dim L=t$, $L\cap K_P=0$, and $\supp L\subseteq\supp L'$.  The reverse
inclusion of feasible classes is immediate.  The two minima are equal.
\uses{prop:associated-pair}%
\proves{thm:relative-weight-recovery}%
\end{proof}

The corresponding relative dimension/length profile is
\[
 K_s(D_P,K_P)=
 \max_{\substack{H\subseteq J\\|H|=s}}
 \bigl(\dim(D_P\cap\F_q^H)-\dim(K_P\cap\F_q^H)\bigr),
\]
where $\F_q^H$ denotes vectors supported on $H$.

\begin{proposition}[Recovery interpretation of the relative profile]
\label{prop:relative-profile}
\coverage{absent}%
For $0\leq s\leq|J|$,
\[
 K_s(D_P,K_P)=
 \max_{\substack{H\subseteq J\\|H|=s}}
 \dim\bigl(W_P\cap\operatorname{span}(G_H)\bigr).
\]
Consequently,
\[
 M_t(D_P,K_P)=\min\{s:K_s(D_P,K_P)\geq t\}.
\]
Thus $K_s$ is the maximum number of independent target linear combinations
recoverable from a fixed budget of $s$ helpers.
\end{proposition}

\begin{proof}
Fix $H$.  Restrict the exact map~\eqref{eq:associated-exact-sequence} to
$D_P\cap\F_q^H$.  Its kernel is $K_P\cap\F_q^H$, and its image is
$W_P\cap\operatorname{span}(G_H)$.  Rank--nullity gives the first equality
after maximizing over $H$.  The second is the standard inverse relation
between the relative dimension/length profile and relative generalized
weights; it also follows directly from Theorem~\ref{thm:relative-weight-recovery}.
\uses{prop:associated-pair, thm:relative-weight-recovery}%
\proves{prop:relative-profile}%
\end{proof}

Standard shortening--puncturing duality gives the complementary failure
interpretation: the relative weights of
$D_P^\perp\subseteq K_P^\perp$ are the minimum numbers of helper failures
that leave prescribed dimensions of $W_P$ ambiguous
\cite{LuoEtAl2005,GeilEtAl2014}.

The strict inequalities
\[
 1\leq M_1(D_P,K_P)<\cdots<M_\ell(D_P,K_P)
\]
therefore say that each additional independent target combination requires a
strictly larger helper union.  These minima solve the single-block problem but
do not yet provide recursive state: the next section restores the functional
labels needed for exact concatenation.  Only when an outer-distance condition
removes every nonzero label does the answer collapse to
$M_t(D_P,K_P)+d(I^\perp)$.

\subsection{Prescribed quotient operations}

The pair construction admits a useful general form.  For finite linear codes
$K\subseteq D\leq\F_q^E$, let $\pi:D\to D/K$ be the quotient map.
Given a finite $\F_q$-space $T$ and a linear map $B:T\to D/K$, put
\[
 \kappa_T^{D/K}(B)=
 \min_{\substack{Y:T\to D\ \mathrm{linear}\\\pi Y=B}}
 |\supp Y(T)|.
\]
Thus a quotient label specifies an operation and a lift specifies its
physical coordinates.  Under $D_P/K_P\cong W_P$, the inclusion of a
target subspace $T$ has cost $\rho_T(I)$, the least helper-union size
for recovering $T$ in $I$: its helper lift has image under
$G_J$ equal to the requested target, and negating the helper coefficients
gives the normalized dual equations.  More generally, minimizing
$\kappa_T^{D/K}(B)$ over injective $B$ with $\dim T=t$ gives $M_t(D,K)$.
Indeed an injective quotient image forces the lift image to be disjoint
from $K$, and every $t$-dimensional subspace disjoint from $K$ supplies
such a map after choosing a basis.

\begin{theorem}[Composition of quotient-labelled lifts]
\label{thm:quotient-lifting-composition}
\coverage{absent}%
Let $E_i$ be disjoint, let $K_i\subseteq D_i\leq\F_q^{E_i}$, and write
$\pi_i:D_i\to L_i=D_i/K_i$.  For a nested pair
$R_0\subseteq R_1\leq\bigoplus_iL_i$, take inverse images inside
$\bigoplus_iD_i$:
\[
 \pi=\bigoplus_i\pi_i,\qquad D=\pi^{-1}(R_1),\qquad K=\pi^{-1}(R_0).
\]
Then $D/K\cong R_1/R_0$.  Writing $\bar\pi:R_1\to R_1/R_0$, every
linear $B:T\to R_1/R_0$ satisfies
\begin{equation}
 \kappa_T^{D/K}(B)=
 \min_{\substack{Z:T\to R_1\ \mathrm{linear}\\\bar\pi Z=B}}
 \sum_i\kappa_T^{D_i/K_i}(\operatorname{pr}_iZ).
 \label{eq:quotient-lifting-composition}
\end{equation}
Minimizing lifts compose to a witness.  Repeated substitution is associative.
\end{theorem}

\begin{proof}
The surjection $D\to R_1/R_0$ induced by $\pi$ has kernel $K$, giving
the asserted isomorphism.  A lift $Y:T\to D$ determines $Z=\pi Y$ with
$\bar\pi Z=B$.  For fixed $Z$, its block maps $Y_i$ are precisely the
independent lifts of $\operatorname{pr}_iZ$.  Their supports are disjoint,
so their minimum costs add.  All quotient maps admit linear sections;
hence these minima are attained and their block witnesses assemble to a
lift of $B$.  This proves both inequalities in the formula.  At further
levels, either order of substitution minimizes over the same compatible
intermediate maps and the same sum of leaf costs.
\proves{thm:quotient-lifting-composition}%
\end{proof}

No common extension field is needed here.  For heterogeneous injective
encoders $e_i:V_i\hookrightarrow\F_q^{E_i}$ and a base-field-linear
outer relation $O\leq\bigoplus_iV_i$, the interface is
$e_i^*:\F_q^{E_i}\to V_i^*$ and compatibility is membership in
$\operatorname{ann}(O)$.  Orthogonality to all encoded outer words proves
this directly.  Target-normalized fibres retain their target constraints
and exclude target coordinates from the cost.  To identify a prescribed
target with the full alphabet $V_j$, the projection $O\to V_j$ must be
surjective; a nonzero projection alone is insufficient.  With nonzero $T$
and at least two blocks, the zero-functional nonconfinement contribution is
$\rho_{j,T}+\min_{i\ne j}d(I_i^\perp)$, by adjoining a minimum dual word
in the cheapest external block.  Extension-field scalar and projective
actions used below are additional structure of that specialization.

\paragraph{Available quotient directions.}
For a coordinate set $A\subseteq E$, the quotient directions realizable
entirely on $A$ form
\[
 L(A)=\pi(D\cap\F_q^A),\qquad
 \dim L(A)=\dim(D\cap\F_q^A)-\dim(K\cap\F_q^A).
\]
Here $\F_q^A$ is embedded by zero-extension.  The formula is rank--nullity
for the restricted quotient map.  A prescribed $B$ has a lift on $A$
exactly when $B(T)\subseteq L(A)$, since the restricted surjection admits
a linear section.  This retains the directions exposed on $A$, as well as
their maximum possible rank; minimizing ordinary words of $D$ would lose
the requirement that the quotient label be nonzero.
